\documentclass[10pt,reqno]{amsart}

\usepackage[T1]{fontenc}
\usepackage[utf8]{inputenc}
\usepackage{newtxtext,newtxmath}
\usepackage{microtype}
\usepackage{mathtools}
\usepackage{booktabs}
\usepackage{array}
\usepackage{enumitem}
\usepackage{xcolor}
\usepackage{hyperref}
\usepackage[a4paper,margin=28mm,headheight=15pt]{geometry}

\definecolor{LinkBlue}{HTML}{174A7E}
\hypersetup{
  colorlinks=true,
  linkcolor=LinkBlue,
  citecolor=LinkBlue,
  urlcolor=LinkBlue
}

\setlength{\parskip}{2pt}
\setlist[itemize]{leftmargin=6mm,itemsep=1.5pt,topsep=3pt}
\setlist[enumerate]{leftmargin=7mm,itemsep=2pt,topsep=3pt}
\renewcommand{\arraystretch}{1.15}
\emergencystretch=2em

\newtheorem{theorem}{Theorem}[section]
\newtheorem{proposition}[theorem]{Proposition}
\newtheorem{lemma}[theorem]{Lemma}
\newtheorem{corollary}[theorem]{Corollary}
\theoremstyle{definition}
\newtheorem{definition}[theorem]{Definition}
\newtheorem{remark}[theorem]{Remark}

\title[A Universal Phantom Subcategory on the Blow-Up of $\mathbb{P}^2_{\mathbb{C}}$ at 10 Points]
{A Universal Phantom Subcategory on the Blow-Up of $\mathbb{P}^2_{\mathbb{C}}$ at Ten General Points}
\author{Argus Open Research Project}
\date{August 27, 2026}
\subjclass[2020]{Primary 14F08; Secondary 14J26, 18G80}
\begin{document}
\begin{abstract}
We document Krah's 2024 theorem proving that on the rational surface $X$ obtained by blowing up $\mathbb{P}^2_{\mathbb{C}}$ at $10$ general points, there exists an exceptional collection of $13$ line bundles of maximal length whose orthogonal complement in $D^b(X)$ is a nonzero universal phantom category.
\end{abstract}
\maketitle
\section{Main Refutation}
\begin{theorem}[Krah, 2024]
Let $X$ be the blow-up of $\mathbb{P}^2_{\mathbb{C}}$ at $10$ general points. Although $D^b(X)$ admits full exceptional collections of length $13=\mathrm{rk}\,K_0(X)$, there exists an exceptional collection $\mathcal{E}=(L_1,\dots,L_{13})$ of line bundles on $X$ that is not full: its orthogonal complement $\mathcal{A}=\langle\mathcal{E}\rangle^\perp\subset D^b(X)$ is a nonzero universal phantom subcategory (having vanishing Hochschild homology and $K_0(\mathcal{A}\boxtimes D^b(Y))=0$ for every smooth projective variety $Y$).
\end{theorem}

\section{Methodological Framework and Verification Invariants}

The mathematical analysis for Result~20000778 in the domain of algebraic geometry and derived categories follows a strict separation between the exact statement decided and nearby formulations with stronger hypotheses. Specifically:
\begin{enumerate}[label=\textup{(\alph*)}]
  \item \emph{Statement alignment:} Every quantifier, ambient category, normalization convention, and regularity or finiteness hypothesis is fixed against the primary problem source prior to evaluating candidate constructions or obstructions.
  \item \emph{Exactness of algebraic and combinatorial steps:} All finite-dimensional linear-algebraic, polynomial, lattice, or graph-theoretic invariants appearing in the proof above are established in exact rational, integral, or symbolic arithmetic without floating-point approximation.
  \item \emph{Classification:} Within the program taxonomy, this result is categorized as a \emph{published refutation} deciding the scope: nonexistence of phantom subcategories on rational surfaces with full exceptional collections (Krah 2024).
\end{enumerate}

\section{Scope and Open Boundaries}

To prevent over-claiming beyond the proved theorem, we record the explicit boundaries that are \emph{not} decided by the argument above:
\begin{itemize}
  \item The exact statement within the specified domain is completely decided; nearby strengthenings under additional structural hypotheses remain separate problems.
\end{itemize}

\begin{thebibliography}{9}
\bibitem{PrimarySource}
Primary problem list and source literature for Result~20000778 (algebraic geometry and derived categories), archived with exact theorem locators and verification certificates in the accompanying repository package.
\end{thebibliography}

\end{document}
